% Nội dung phần 4: Nhận xét và kết luận.

\section{Nhận xét và kết luận}

\subsection{Nhận xét về kết quả phân tích}

\subsubsection{Mô hình nào hoạt động tốt nhất và tại sao}

Theo quy tắc lựa chọn được chốt trước khi mở dữ liệu đánh giá,
\textbf{Logistic GLM với tầng predictor history-enhanced} là mô hình được chọn.
Lý do không nằm ở khả năng phân biệt -- về mặt này ba mô hình gần như không phân
biệt được, với ROC--AUC trên temporal test lần lượt là 0,662 (LPM), 0,666
(Logistic GLM) và 0,666 (Binomial GAM), các khoảng tin cậy bootstrap chồng lấn
gần hoàn toàn (Bảng~\ref{tab:final-metrics}).

Sự khác biệt thực chất nằm ở ba tiêu chí còn lại:

\begin{itemize}
  \item \textbf{Tính hợp lệ của xác suất.} LPM tạo ra 1,706\% dự đoán nằm ngoài
  khoảng $[0,1]$ trên temporal test -- những giá trị không thể diễn giải như xác
  suất. Logistic GLM và Binomial GAM theo cấu trúc không bao giờ mắc lỗi này.

  \item \textbf{Calibration.} Calibration slope của LPM chỉ đạt 0,868, trong khi
  Logistic GLM đạt 0,934 và GAM đạt 0,942 (lý tưởng là 1). Xác suất do LPM sinh
  ra cực đoan hơn mức dữ liệu thực sự cho phép.

  \item \textbf{Khả năng diễn giải và độ phức tạp.} Logistic GLM và Binomial GAM
  cho log loss (0,379) và Brier (0,113) giống nhau đến ba chữ số thập phân, tức
  GAM \emph{không} mua thêm được gì bằng độ phức tạp của nó. Thêm vào đó,
  concurvity giữa các term của GAM lên tới 0,994 (Bảng~\ref{tab:collinearity}),
  khiến các hàm trơn của nó không thể diễn giải riêng lẻ. Ngược lại, GVIF điều
  chỉnh lớn nhất của Logistic GLM chỉ 2,424, nên các hệ số của nó đọc được.
\end{itemize}

Sự tương đương giữa GAM và mô hình tuyến tính có một lý giải rõ ràng từ dữ liệu:
sau phép biến đổi $\log_2(\texttt{order})$, quan hệ giữa tiến trình session và
xác suất kết thúc đã gần như tuyến tính. Hình~\ref{fig:order-probability} cho
thấy ba đường dự đoán trùng nhau trong khoảng \texttt{order} từ 1 đến 20, và
Bảng~\ref{tab:kcheck} cho thấy các hàm trơn của biến giá chỉ có
$\text{edf} \approx 1$--$2$, tức gần như là đường thẳng. Nói cách khác, phần phi
tuyến mà GAM có thể khai thác trong bộ dữ liệu này là rất nhỏ.

\subsubsection{Các biến có ảnh hưởng lớn nhất}

Dựa trên kiểm định Wald robust theo cụm session (Bảng~\ref{tab:terms}), năm nhóm
biến có bằng chứng liên hệ mạnh nhất với \texttt{Exit} là
$\log_2(\text{Order})$, số danh mục duy nhất đã xem, hành vi quay lại sản phẩm
cũ, \texttt{Page} và \texttt{Country}. Ba hiệu ứng đáng chú ý nhất về độ lớn:

\begin{itemize}
  \item Click vào một sản phẩm \emph{đã xem trước đó} trong cùng session làm odds
  kết thúc tăng gấp 2,218 lần (CI 95\%: 2,064--2,384) -- hiệu ứng đơn biến mạnh
  nhất trong mô hình.
  \item Mỗi danh mục mới được xem thêm làm odds kết thúc nhân với 1,562
  (CI 95\%: 1,509--1,617).
  \item Ngược chiều, mỗi lần \texttt{order} tăng gấp đôi làm odds kết thúc nhân
  với 0,640 (CI 95\%: 0,623--0,657): session càng đi sâu thì càng ít khả năng
  dừng ở click kế tiếp.
\end{itemize}

Một phát hiện có giá trị phương pháp luận là \textbf{nhóm biến lịch sử duyệt quan
trọng hơn hẳn đặc điểm sản phẩm}. Thêm toàn bộ giá, danh mục, màu sắc, vị trí
ảnh, kiểu ảnh và trang hiển thị chỉ nâng ROC--AUC từ 0,589 lên 0,611; trong khi
nhóm biến lịch sử nâng tiếp lên 0,662 -- đóng góp hơn gấp đôi về ROC--AUC và gấp
ba lần về log loss. Cả hai bước
đều có robust joint $p < 0{,}001$, nhưng ý nghĩa thống kê và độ lớn thực tế là
hai chuyện khác nhau, và ở đây độ lớn nghiêng hẳn về nhóm biến hành vi.

Riêng \texttt{Price} tuy đạt ý nghĩa thống kê ($p = 0{,}008$) nhưng hiệu ứng rất
nhỏ (OR 1,032 cho mỗi 10 USD). Với cỡ mẫu 116.095 click của tập train, ngay cả
những liên hệ rất yếu cũng dễ đạt $p < 0{,}05$; vì vậy nhóm nhất quán đọc p-value
cùng với độ lớn hiệu ứng và khoảng tin cậy.

\subsubsection{Kết luận về kiểm định hai nhóm giá}

Ở mức ý nghĩa $\alpha = 0{,}05$, nhóm \textbf{bác bỏ giả thuyết $H_0$}: xác suất
\texttt{Exit} của hai nhóm giá không bằng nhau. Bằng chứng đến từ nhiều nguồn
nhất quán:

\begin{itemize}
  \item Kiểm định hai tỷ lệ thô: $\chi^2 = 10{,}08$, $p = 0{,}002$.
  \item Adjusted odds ratio với sai số chuẩn cluster-robust theo session:
  1,071 (CI 95\%: 1,002--1,144), $p = 0{,}043$.
  \item Chênh lệch thô 1,19 điểm phần trăm với cluster-bootstrap CI 95\% từ 0,46
  đến 1,93 điểm phần trăm -- không chứa 0.
  \item Cùng dấu và độ lớn tương đương trên cả ba giai đoạn thời gian
  (Bảng~\ref{tab:ab-temporal}).
\end{itemize}

Tuy nhiên, kết luận này cần ba dè dặt quan trọng. \emph{Thứ nhất}, độ lớn hiệu
ứng rất nhỏ: adjusted risk difference chỉ 0,82 điểm phần trăm với cận dưới khoảng
tin cậy là 0,03 -- gần như chạm 0. \emph{Thứ hai}, hiệu ứng không đồng đều giữa
các danh mục: sau hiệu chỉnh Holm chỉ riêng Trousers còn có ý nghĩa (2,32 điểm
phần trăm, $p = 0{,}008$), còn Skirts, Blouses và Sale đều có khoảng tin cậy chứa
0, với kiểm định tương tác nằm ngay ranh giới ($p = 0{,}051$). \emph{Thứ ba} và
quan trọng nhất, balance diagnostics (Hình~\ref{fig:ab-balance}) xác nhận hai
nhóm không tương đương như trong một thí nghiệm ngẫu nhiên. Đây là dữ liệu quan
sát, nên kết quả chỉ chứng minh một \textbf{mối liên hệ}, hoàn toàn không chứng
minh rằng việc tăng giá \emph{gây ra} việc kết thúc session.

\subsection{Kết luận tổng thể và khuyến nghị}

\subsubsection{Những hiểu biết rút ra}

\textbf{1. Rủi ro rời bỏ tập trung ở đầu session.} Xác suất kết thúc giảm từ
20,47\% tại click đầu tiên xuống 7,64\% sau click thứ 20 -- giảm hơn một nửa
(Bảng~\ref{tab:hazard}). Người dùng vượt qua được vài click đầu có xu hướng ở lại
lâu hơn nhiều.

\textbf{2. Hành vi trong phiên nói nhiều hơn thuộc tính sản phẩm.} Biết người
dùng \emph{đã làm gì} trong session (đã xem bao nhiêu danh mục, có quay lại sản
phẩm cũ không, có nhảy trang không) dự báo tốt hơn hẳn việc biết họ \emph{đang
xem sản phẩm nào}.

\textbf{3. Quay lại sản phẩm cũ là tín hiệu cảnh báo mạnh nhất.} Odds kết thúc
tăng gấp 2,218 lần. Cần lưu ý đây là quan hệ hai chiều về mặt diễn giải: hành vi
này vừa có thể là dấu hiệu người dùng đã tìm được thứ cần và sắp rời đi, vừa có
thể là dấu hiệu họ đang loanh quanh không tìm thấy gì mới. Dữ liệu hiện có không
phân biệt được hai kịch bản này vì thiếu thông tin giao dịch.

\textbf{4. Nhưng khả năng dự báo tổng thể còn hạn chế.} Đây là kết luận trung
thực cần nêu rõ: ROC--AUC 0,666 chỉ ở mức trung bình, Brier Skill Score 0,039
nghĩa là mô hình chỉ tốt hơn dự báo hằng khoảng 4\%, và tại threshold tối ưu
$F_1$ thì Precision chỉ 20,94\% -- cứ 5 cảnh báo thì khoảng 1 cảnh báo đúng.
Nguyên nhân chính là dữ liệu thiếu những tín hiệu quan trọng nhất: thời lượng
dừng lại trên mỗi trang, thao tác thêm vào giỏ hàng, và lịch sử người dùng qua
nhiều phiên.

\subsubsection{Khuyến nghị hành động}

\textbf{Về mặt vận hành:}

\begin{enumerate}
  \item \textbf{Ưu tiên can thiệp ở vài click đầu tiên.} Vì gần một phần năm
  session kết thúc ngay sau click đầu, các nỗ lực cải thiện trải nghiệm nên tập
  trung vào trang đầu tiên người dùng nhìn thấy. Lưu ý \texttt{Order} là một chỉ
  báo về giai đoạn session, không phải một biến có thể can thiệp trực tiếp.

  \item \textbf{Tính các đặc trưng hành vi theo thời gian thực.} Số sản phẩm và
  danh mục đã xem, cờ quay lại sản phẩm cũ, cờ chuyển danh mục -- toàn bộ đều
  tính được ngay tại click hiện tại mà không cần biết tương lai của session, nên
  triển khai được trong hệ thống trực tuyến.

  \item \textbf{Không tự động hóa quyết định dựa trên mô hình hiện tại.} Với
  Precision khoảng 21\%, việc tự động kích hoạt ưu đãi cho mọi phiên bị dự báo
  ``sắp rời'' sẽ lãng phí phần lớn ngân sách. Nếu dùng, chỉ nên dùng như một tín
  hiệu xếp hạng ưu tiên trong một hệ thống có chi phí can thiệp thấp.

  \item \textbf{Theo dõi calibration khi triển khai, không chỉ AUC.} Calibration
  intercept trên temporal test là $-0{,}127$, âm hơn so với mức $-0{,}059$ trên
  validation, cho thấy chất lượng xác suất suy giảm theo thời gian ngay cả khi
  AUC giữ nguyên. Mô hình cần được hiệu chỉnh lại định kỳ.
\end{enumerate}

\textbf{Về mặt thu thập dữ liệu:} giá trị gia tăng lớn nhất không đến từ mô hình
phức tạp hơn mà đến từ dữ liệu tốt hơn. Ba trường cần bổ sung, theo thứ tự ưu
tiên, là: dấu thời gian đầy đủ của từng click (để tính thời lượng dừng), sự kiện
thêm vào giỏ hàng và mua hàng (để chuyển từ dự báo ``rời phiên'' sang dự báo
``chuyển đổi''), và định danh người dùng xuyên phiên.

\textbf{Về nhóm giá:} không nên thay đổi chính sách giá dựa trên kết quả của báo
cáo này. Bằng chứng thu được là liên hệ trên dữ liệu quan sát, với độ lớn nhỏ và
không đồng đều giữa các danh mục. Bước đi đúng đắn là chạy một thí nghiệm ngẫu
nhiên thực sự theo thiết kế được trình bày ở Mục~\ref{subsec:future-ab}.
